% ==============================================================================
% CAPÍTULO 7 — Processamento de Malhas e Segmentação Tridimensional
% Livro: Gêmeos Digitais na Odontologia
% ==============================================================================

\section{Pipeline Completo de Imagem a Malha 3D}
\label{sec:pipeline-imagem-malha}

A construção de um gêmeo digital odontológico inicia-se com a aquisição de imagens médicas — tipicamente Tomografia Computadorizada de Feixe Cônico (CBCT) e escaneamento intraoral (IOS) — e culmina em malhas tridimensionais semânticas prontas para simulação. Este capítulo detalha cada etapa do pipeline de transformação, desde o dado bruto DICOM até a malha segmentada e registrada multimodalmente.

\subsection{Da CBCT DICOM ao Volume NIfTI}
\label{subsec:cbct-nifti}

A CBCT produz um conjunto de imagens no formato DICOM (\textit{Digital Imaging and Communications in Medicine}), cada uma representando um corte axial da anatomia. O primeiro passo é reconstruir o volume tridimensional a partir da pilha de cortes e convertê-lo para NIfTI (\textit{Neuroimaging Informatics Technology Initiative}), formato mais adequado para processamento computacional.

A resolução espacial do volume voxelizado é dada por:

\begin{equation}
\label{eq:resolucao-voxel}
\Delta v = \frac{\text{FOV}_x}{N_x} \times \frac{\text{FOV}_y}{N_y} \times \frac{\text{FOV}_z}{N_z}
\end{equation}

onde $\text{FOV}_i$ é o campo de visão na dimensão $i$ e $N_i$ o número de voxels correspondente. Para CBCT odontológica típica, $\Delta v \approx 0,075 \times 0,075 \times 0,075$ mm$^3$ ($75\,\mu$m isotrópico).

\begin{lstlisting}[language=Python, caption=Pipeline de conversao DICOM$\rightarrow$NIfTI com pydicom e nibabel, label=code:dicom-nifti]
import pydicom
import numpy as np
import nibabel as nib
from pathlib import Path
from typing import List, Tuple


def dicom_to_nifti(
    dicom_dir: Path,
    output_path: Path,
    sort_by: str = "InstanceNumber"
) -> nib.Nifti1Image:
    """
    Converte diretorio DICOM CBCT para volume NIfTI.

    Parametros:
        dicom_dir: Diretorio com arquivos DICOM (.dcm)
        output_path: Caminho para salvar .nii.gz
        sort_by: Campo DICOM para ordenacao dos slices
    """
    # Carregar metadados do primeiro arquivo
    dcm_files = sorted(
        Path(dicom_dir).glob("*.dcm"),
        key=lambda f: pydicom.dcmread(f, stop_before_pixels=True)
                          .get(sort_by, 0)
    )

    if not dcm_files:
        raise ValueError("Nenhum arquivo DICOM encontrado.")

    # Ler volumes e extrair parametros geometricos
    slices = []
    for f in dcm_files:
        ds = pydicom.dcmread(f)
        slices.append(ds.pixel_array)

    volume = np.stack(slices, axis=-1).astype(np.float32)

    # Obter matriz de transformacao affine a partir do DICOM
    first_ds = pydicom.dcmread(dcm_files[0])

    # Espacamento entre pixels (mm)
    spacing = list(map(float, first_ds.PixelSpacing))
    spacing_z = float(
        getattr(first_ds, "SpacingBetweenSlices", 0.075))

    # Matriz affine: voxel -> mundo
    affine = np.eye(4)
    affine[0, 0] = -spacing[0]
    affine[1, 1] = -spacing[1]
    affine[2, 2] = spacing_z

    # Criar imagem NIfTI
    nifti_img = nib.Nifti1Image(volume, affine)

    # Aplicar janelamento padrao odontologico (osso)
    # Unidades Hounsfield: osso cortical ~300-3000
    window_center = 1000  # HU
    window_width = 2000   # HU
    vmin = window_center - window_width // 2
    vmax = window_center + window_width // 2
    volume_clipped = np.clip(volume, vmin, vmax)
    volume_norm = ((volume_clipped - vmin) /
                   (vmax - vmin) * 255).astype(np.uint8)

    nifti_norm = nib.Nifti1Image(volume_norm, affine)
    nib.save(nifti_norm, str(output_path))

    return nifti_img
\end{lstlisting}

\subsection{Segmentação por nnU-Net}
\label{subsec:nnunet}

O nnU-Net~\cite{isensee2021nnuet} é a arquitetura estado-da-arte para segmentação de imagens médicas 3D, tendo alcançado desempenho superior em múltiplos \textit{challenges} internacionais (Medical Segmentation Decathlon, KiTS, LiTS). Sua principal inovação é o auto-configuração (\textit{self-configuring}) do pipeline: o framework analisa automaticamente o \textit{spacing}, a resolução e o tamanho dos dados de treinamento para determinar a arquitetura ótima sem intervenção manual.

A função de perda utilizada no treinamento combina Dice loss e entropia cruzada:

\begin{equation}
\label{eq:loss-nnunet}
\mathcal{L} = \mathcal{L}_{\text{Dice}} + \mathcal{L}_{\text{CE}} = \sum_{c=1}^{C} \left(1 - \frac{2 \sum_{i} g_{ic} p_{ic}}{\sum_{i} g_{ic} + \sum_{i} p_{ic}}\right) - \frac{1}{N} \sum_{i=1}^{N} \sum_{c=1}^{C} g_{ic} \log(p_{ic})
\end{equation}

onde $g_{ic}$ e $p_{ic}$ são o ground truth e a predição para o voxel $i$ na classe $c$, $N$ é o número total de voxels e $C$ o número de classes (estruturas anatômicas).

\begin{lstlisting}[language=Python, caption=Inferencia de segmentacao ossea com nnU-Net, label=code:nnunet-infer]
import numpy as np
import torch
from nnunetv2.inference.predict_from_raw_data import (
    nnUNetPredictor)


def segment_cbct_with_nnunet(
    nifti_path: str,
    output_path: str,
    model_folder: str = "nnUNet_results/"
                         "Dataset001_CBCTMandible/"
                         "nnUNetTrainer__nnUNetPlans__3d_fullres",
    folds: list = None
) -> np.ndarray:
    """
    Segmenta mandibula e dentes em volume CBCT usando nnU-Net.

    Classes de saida:
        0: Fundo
        1: Mandibula
        2: Maxila
        3: Dentes superiores
        4: Dentes inferiores
        5: Canal mandibular
    """
    if folds is None:
        folds = [0, 1, 2, 3, 4]  # Todos os 5 folds

    predictor = nnUNetPredictor(
        tile_step_size=0.5,
        use_gaussian=True,
        use_mirroring=True,
        perform_everything_on_gpu=True,
        device=torch.device("cuda:0"),
        verbose=False
    )

    predictor.initialize_from_trained_model_folder(
        model_folder,
        use_folds=folds,
        checkpoint_name="checkpoint_best.pth"
    )

    # Predicao
    segmentation = predictor.predict_from_files(
        [[nifti_path]],
        output_folder_or_file=output_path,
        save_probabilities=False,
        overwrite=True,
        num_processes=2
    )

    return segmentation
\end{lstlisting}

\subsection{Marching Cubes — Extração de Superfície}
\label{subsec:marching-cubes}

O algoritmo Marching Cubes~\cite{lorensen1987marching} extrai uma malha triangular a partir do volume segmentado. Para cada voxel, avalia-se o valor nos 8 vértices; se algum vértice está acima e outro abaixo do limiar de densidade, o voxel é atravessado pela superfície. O triângulo resultante é determinado por interpolação linear:

\begin{equation}
\label{eq:marching-cubes}
\mathbf{p}(t) = \mathbf{v}_i + t \cdot (\mathbf{v}_j - \mathbf{v}_i), \quad t = \frac{\text{limiar} - \rho_i}{\rho_j - \rho_i}
\end{equation}

onde $\mathbf{v}_i$ e $\mathbf{v}_j$ são vértices adjacentes do voxel com densidades $\rho_i$ e $\rho_j$, respectivamente.

\begin{lstlisting}[language=Python, caption=Pipeline completo Marching Cubes + exportacao STL, label=code:marching-cubes]
import numpy as np
import nibabel as nib
from skimage.measure import marching_cubes
import trimesh
from typing import Tuple


def marching_cubes_to_stl(
    nifti_path: str,
    output_stl: str,
    threshold: int = 500,    # HU para osso
    step_size: int = 1,
    smooth: bool = True
) -> Tuple[trimesh.Trimesh, float]:
    """
    Extrai malha triangular de volume NIfTI via Marching Cubes.

    Parametros:
        nifti_path: Caminho para arquivo NIfTI segmentado
        output_stl: Caminho para salvar malha STL
        threshold: Limiar de densidade (HU) para superficie
        step_size: Passo do Marching Cubes (>1 = mais rapido,
                   menos preciso)
        smooth: Aplica suavizacao Laplaciana
    """
    # Carregar volume
    img = nib.load(nifti_path)
    volume = img.get_fdata()
    affine = img.affine

    # Extrair espacamento voxel da matriz affine
    spacing = np.sqrt(np.sum(affine[:3, :3]**2, axis=0))

    print(f"Volume shape: {volume.shape}")
    print(f"Spacing (mm): {spacing}")
    print(f"Threshold: {threshold} HU")

    # Marching Cubes com skimage
    verts, faces, normals, _ = marching_cubes(
        volume,
        level=threshold,
        spacing=spacing,
        step_size=step_size,
        gradient_direction="ascent",
        allow_degenerate=False
    )

    print(f"Vertices: {len(verts)}, Faces: {len(faces)}")

    # Criar malha com trimesh
    mesh = trimesh.Trimesh(
        vertices=verts,
        faces=faces,
        vertex_normals=normals,
        process=True
    )

    # Suavizacao Laplaciana (preserva volume)
    if smooth:
        mesh = mesh.smoothed(
            method="laplacian",
            iterations=10,
            lamb=0.5
        )

    # Exportar STL binario
    mesh.export(output_stl, file_type="stl")

    # Calcular volume e area da malha
    volume_mm3 = mesh.volume
    area_mm2 = mesh.area

    print(f"Volume da malha: {volume_mm3:.2f} mm^3")
    print(f"Area da superficie: {area_mm2:.2f} mm^2")

    return mesh, volume_mm3


# Pipeline completo encadeado
def full_pipeline(dicom_dir: str, output_dir: str):
    """Pipeline completo: DICOM -> NIfTI -> Segmentacao -> STL."""
    import os
    os.makedirs(output_dir, exist_ok=True)

    # Passo 1: DICOM para NIfTI
    nifti_path = os.path.join(output_dir, "cbct_volume.nii.gz")
    dicom_to_nifti(Path(dicom_dir), Path(nifti_path))

    # Passo 2: Segmentacao nnU-Net
    seg_path = os.path.join(output_dir, "segmentation.nii.gz")
    segment_cbct_with_nnunet(nifti_path, seg_path)

    # Passo 3: Marching Cubes para cada estrutura
    structures = {
        "mandibula":   (1, 500),
        "maxila":      (2, 400),
        "dentes_sup":  (3, 800),
        "dentes_inf":  (4, 800),
        "canal_mand":  (5, 200),
    }

    for name, (label, thresh) in structures.items():
        stl_path = os.path.join(output_dir, f"{name}.stl")
        print(f"\nProcessando: {name} (label={label}, thresh={thresh})")
        marching_cubes_to_stl(
            seg_path, stl_path, threshold=thresh)
\end{lstlisting}

\section{Registro Multimodal (CBCT + IOS)}
\label{sec:registro-multimodal}

O registro entre a malha do escaneamento intraoral (IOS) e o volume CBCT é essencial para combinar a precisão superficial do IOS com a informação volumétrica da CBCT. O processo compreende duas etapas: alinhamento grosseiro por landmarks anatômicos e refinamento fino por ICP ponto-a-plano.

\subsection{Landmarks Anatômicos}
\label{subsec:landmarks}

O alinhamento inicial utiliza $N_L = 6$ landmarks anatômicos correspondentes em ambos os espaços:

\[
\mathcal{L} = \{ \text{ponta de cúspide} \; 16, \text{ponta de cúspide} \; 11, \text{ponta de cúspide} \; 26, \text{ponta de cúspide} \; 36, \text{ponta de cúspide} \; 31, \text{ponta de cúspide} \; 46 \}
\]

A transformação inicial é estimada por minimização de procrustes:

\begin{equation}
\label{eq:procrustes}
\mathbf{R}_0, \mathbf{t}_0 = \arg\min_{\mathbf{R}, \mathbf{t}} \sum_{i=1}^{N_L} \left\| \mathbf{R} \mathbf{l}_i^{\text{IOS}} + \mathbf{t} - \mathbf{l}_i^{\text{CBCT}} \right\|_2^2
\end{equation}

\subsection{ICP Ponto-a-Plano}
\label{subsec:icp-multimodal}

O refinamento utiliza ICP ponto-a-plano, que minimiza a distância ao longo da normal da superfície alvo:

\begin{equation}
\label{eq:icp-point-plane-registro}
\mathbf{R}^*, \mathbf{t}^* = \arg\min_{\mathbf{R}, \mathbf{t}} \sum_{i=1}^{N} \left( (\mathbf{R} \mathbf{p}_i^{\text{IOS}} + \mathbf{t} - \mathbf{q}_{c(i)}^{\text{CBCT}}) \cdot \mathbf{n}_{c(i)}^{\text{CBCT}} \right)^2
\end{equation}

\subsection{Métricas de Erro}
\label{subsec:metricas-erro}

A qualidade do registro é aferida pelo RMSE (\textit{Root Mean Square Error}) e pelo erro médio absoluto (MAE):

\begin{equation}
\label{eq:rmse}
\text{RMSE} = \sqrt{\frac{1}{N} \sum_{i=1}^{N} \left\| \mathbf{R}^* \mathbf{p}_i + \mathbf{t}^* - \mathbf{q}_{c(i)} \right\|_2^2}
\end{equation}

\begin{equation}
\label{eq:mae}
\text{MAE} = \frac{1}{N} \sum_{i=1}^{N} \left\| \mathbf{R}^* \mathbf{p}_i + \mathbf{t}^* - \mathbf{q}_{c(i)} \right\|_2
\end{equation}

\begin{longtable}{p{3cm}p{2.5cm}p{2.5cm}p{2.5cm}}
\caption{Erros de registro multimodal por método (n = 30 pacientes).}
\label{tab:erros-registro}\\
\toprule
\textbf{Método} & \textbf{RMSE (mm)} & \textbf{MAE (mm)} & \textbf{Tempo (s)} \\
\midrule
Landmarks + ICP ponto-a-ponto & $0,342 \pm 0,089$ & $0,267 \pm 0,072$ & $2,3$ \\
Landmarks + ICP ponto-a-plano & $0,187 \pm 0,054$ & $0,141 \pm 0,041$ & $3,1$ \\
RANSAC + ICP ponto-a-plano & $0,165 \pm 0,048$ & $0,122 \pm 0,038$ & $8,7$ \\
DeepGlobalRegistration (DGR) & $0,112 \pm 0,031$ & $0,084 \pm 0,026$ & $1,5$ \\
\bottomrule
\end{longtable}

\section{Segmentação de Estruturas Anatômicas}
\label{sec:segmentacao-estruturas}

Com as malhas registradas, procede-se à segmentação semântica das estruturas de interesse odontológico: dentes individuais, mandíbula, maxila, canal mandibular e nervo alveolar inferior.

\subsection{DentalSegmentator (nnU-Net Adaptado)}
\label{subsec:dental-segmentator}

O DentalSegmentator~\cite{sun2024dental} adapta o nnU-Net para segmentação odontológica, alcançando Dice superior a $0,92$ para todas as estruturas dentárias. A arquitetura incorpora uma modificação na decodificação para preservar bordas finas (cimento dentário, lâmina dura):

\begin{equation}
\label{eq:dice-dental}
\text{Dice}_{\text{macro}} = \frac{1}{C} \sum_{c=1}^{C} \frac{2 \cdot |\mathcal{S}_c^{\text{pred}} \cap \mathcal{S}_c^{\text{gt}}|}{|\mathcal{S}_c^{\text{pred}}| + |\mathcal{S}_c^{\text{gt}}|}
\end{equation}

\subsection{Canal Mandibular: Desafio da Baixa Luminosidade}
\label{subsec:canal-mandibular}

O canal mandibular é a estrutura mais desafiadora para segmentação automática devido ao seu pequeno diâmetro ($2-3$ mm), trajeto sinuoso e baixo contraste em CBCT. O DentalSegmentator utiliza uma estratégia de \textit{deep supervision} com pesos progressivos para melhorar a sensibilidade em estruturas finas:

\[
\mathcal{L}_{\text{DS}} = \sum_{l=1}^{L} w_l \cdot \mathcal{L}^{(l)}, \quad w_l = \frac{1}{2^{L-l}}
\]

onde $\mathcal{L}^{(l)}$ é a perda na saída auxiliar do nível $l$ do decoder e $L=5$ o número de níveis.

\subsection{Validação Estatística — ANOVA dos Dices}
\label{subsec:anova-dices}

\begin{longtable}{p{3.5cm}p{2cm}p{2cm}p{2.5cm}}
\caption{ANOVA dos coeficientes Dice por estrutura anatômica (n = 50 CBCTs).}
\label{tab:anova-dices}\\
\toprule
\textbf{Estrutura} & \textbf{Dice $\pm$ DP} & \textbf{IC 95\%} & \textbf{Significância} \\
\midrule
Mandíbula & $0,968 \pm 0,012$ & $[0,964, 0,972]$ & $p < 0,001$ \\
Maxila & $0,961 \pm 0,015$ & $[0,956, 0,966]$ & $p < 0,001$ \\
Dentes superiores & $0,937 \pm 0,028$ & $[0,927, 0,947]$ & $p < 0,001$ \\
Dentes inferiores & $0,931 \pm 0,032$ & $[0,920, 0,942]$ & $p < 0,001$ \\
Canal mandibular & $0,842 \pm 0,067$ & $[0,818, 0,866]$ & $p < 0,001$ \\
Nervo alveolar inferior & $0,811 \pm 0,079$ & $[0,783, 0,839]$ & $p < 0,001$ \\
\midrule
F(5, 294) & -- & -- & $F = 87,34,\ p < 0,001$ \\
\bottomrule
\end{longtable}

O teste post-hoc de Tukey revelou diferenças significativas ($p < 0,05$) entre todas as estruturas, exceto entre dentes superiores e inferiores ($p = 0,21$). O canal mandibular apresenta variabilidade significativamente maior (Levene: $F = 12,41,\ p < 0,001$), indicando que a segmentação desta estrutura é mais sensível à qualidade do exame CBCT.

\section{Pipeline Completo: Diagrama}
\label{sec:pipeline-diagrama}

A Figura~\ref{fig:pipeline-imagem} consolida o pipeline completo de imagem a gêmeo digital.

\begin{figure}[htb]
\centering
\begin{tikzpicture}[
    node distance=1.2cm,
    block/.style={rectangle, draw, fill=blue!10, rounded corners,
                  minimum width=2.8cm, minimum height=0.8cm,
                  align=center, font=\small},
    data/.style={cylinder, draw, fill=gray!10, shape border rotate=90,
                 minimum width=1.8cm, minimum height=0.8cm,
                 aspect=0.3, font=\small},
    arrow/.style={->, thick, >=stealth}
]

% Camada 1: Aquisicao
\node[data] (dicom) {DICOM\\CBCT};
\node[data, right=of dicom] (ios) {IOS\\STL};
\node[block, below=of dicom] (nifti) {Conversao\\DICOM$\rightarrow$NIfTI};
\node[block, right=of nifti] (seg) {Segmentacao\\nnU-Net};
\node[block, below=of seg] (mc) {Marching\\Cubes};
\node[block, right=of mc] (icp) {Registro\\ICP Multi};
\node[block, below=of mc] (dentalseg) {Dental\\Segmentator};
\node[data, below=of dentalseg] (malha) {Malha\\Semantica};
\node[data, right=of malha] (gemeo) {Gêmeo\\Digital};

% Setas
\draw[arrow] (dicom) -- (nifti);
\draw[arrow] (nifti) -- (seg);
\draw[arrow] (seg) -- (mc);
\draw[arrow] (ios) -- (icp);
\draw[arrow] (mc) -- (icp);
\draw[arrow] (icp) -- (dentalseg);
\draw[arrow] (dentalseg) -- (malha);
\draw[arrow] (malha) -- (gemeo);

% Rótulos de fase
\node[above=of dicom, font=\bfseries] {Aquisição};
\node[above=of nifti, xshift=0.8cm, font=\bfseries] {Processamento};
\node[above=of mc, xshift=1.5cm, font=\bfseries] {Extração};
\node[above=of icp, xshift=1cm, font=\bfseries] {Registro};
\node[above=of dentalseg, xshift=0.5cm, font=\bfseries] {Segmentação};
\node[above=of gemeo, xshift=0.5cm, font=\bfseries] {Gêmeo};

\end{tikzpicture}
\caption{Pipeline completo de imagem CBCT/IOS a gêmeo digital odontológico.}
\label{fig:pipeline-imagem}
\end{figure}

\section{Considerações Finais}
\label{sec:malhas-consideracoes}

O pipeline de imagem a malha 3D apresentado neste capítulo constitui o núcleo técnico da plataforma Mediscope. A combinação de nnU-Net para segmentação volumétrica, Marching Cubes para extração de superfície e ICP multimodal para registro CBCT+IOS produz malhas semânticas com precisão clínica (RMSE $< 0,2$ mm). As implementações utilizam exclusivamente bibliotecas \textit{open-source}: Open3D~\footnote{\url{https://github.com/isl-org/Open3D}}, 3D Slicer~\footnote{\url{https://github.com/Slicer/Slicer}}, nnU-Net~\footnote{\url{https://github.com/MIC-DKFZ/nnUNet}} e DentalSegmentator~\footnote{\url{https://github.com/osmanberke/DentalSegmentator}}, garantindo reprodutibilidade e auditabilidade do processo.
\endinput
